\section{Component growth and a boundary under grouping}
\label{sec:components}

For complex vector spaces $V_1,\ldots,V_k$ and a real $a\geq0$, write
\[
 Z_a(V_1,\ldots,V_k)
   =\{T\in V_1\otimes\cdots\otimes V_k:\AR(T)\leq a\}.
\]
These are the full sublevels, which are Zariski closed by
\cite[Theorem~1.2]{CHNVZ2024}. Ranks and asymptotic ranks always refer
to the displayed factors. Grouping a collection of factors for a matrix
flattening does not change this convention.

\subsection{The local rank-one chart in general dimension}
\label{sec:components-chart}

Let $A,B,C$ have the same dimension $n\geq3$, put $W=A\otimes B\otimes C$,
and let $D$ have dimension $m\geq2$. For $T\in W$, set
\[
 E_{AB}(T)=\im(C^*\longrightarrow A\otimes B),
\]
and define $E_{AC}(T)$ and $E_{BC}(T)$ in the same way. Using the
natural permutations of tensor factors, define
\begin{equation}\label{eq:components-K}
 K(T)=(E_{AB}(T)\otimes C)\cap(E_{AC}(T)\otimes B)
                    \cap(A\otimes E_{BC}(T))\subseteq W.
\end{equation}
Thus $T\in K(T)$, and $T$ is concise precisely when the three
pair-image spaces have dimension $n$.

The appended-minor construction of Section~\ref{sec:geometry} extends
without change: to each $n^2$-by-$n$ flattening of $T$, append each
column of the corresponding flattening of a variable $S$, and take
all $(n+1)$-row determinants. Their coefficients in $S$ form a
polynomial matrix $L_n(T)$ with $n^3$ columns. Repeated columns give
$L_n(T)T=0$ for every $T$, and on the concise locus one has
\[
 \ker L_n(T)=K(T).
\]
Indeed, a full-rank $n$-column matrix acquires no additional rank
on appending a column exactly when that column lies in its image.
Consequently
\begin{equation}\label{eq:components-Omega}
 \Omega_n=\{T\in W:T\text{ is concise},\ \rank L_n(T)\geq n^3-1\}
\end{equation}
is Zariski open, and $K(T)=\C T$ for every $T\in\Omega_n$.

Choose a basis $d_0,\ldots,d_{m-1}$ of $D$ and write
$F=\sum_{j=0}^{m-1}T_j\otimes d_j$. Define
\[
 \begin{split}
 \mathcal F_n={}&\{F:\rank F_{AB\mid CD},\ \rank F_{AC\mid BD},
                          \ \rank F_{BC\mid AD}\leq n\},\\
 X={}&\{F:\rank F_{ABC\mid D}\leq1\},\qquad
 U=\{F:T_0\in\Omega_n\}.
 \end{split}
\]
Here the first line requires each of the three ranks to be at most $n$.

\begin{proposition}\label{prop:components-chart}
The full-ambient open set $U$ satisfies
\begin{equation}\label{eq:components-chart-equality}
 \mathcal F_n\cap U=X\cap U.
\end{equation}
If $\Omega_n\ne\varnothing$, then $X$ is an irreducible component of
$\mathcal F_n$ of dimension $n^3+m-1$. In the four-factor rank convention
$(A,B,C,D)$, for every $a\geq0$ one has
\begin{equation}\label{eq:components-sublevel-X}
 Z_a(A,B,C,D)\cap X
   =\{T\otimes v:T\in Z_a(A,B,C),\ v\in D\}.
\end{equation}
For $0\leq a<n+1$, its intersection with $U$ is isomorphic to
\begin{equation}\label{eq:components-sublevel-chart}
 \bigl(Z_a(A,B,C)\cap\Omega_n\bigr)\times\mathbb A^{m-1}.
\end{equation}
\end{proposition}

\begin{proof}
Suppose $F\in U$. The $AB\mid CD$ flattening concatenates the columns
of all the $T_j$, while the first slice already has an $n$-dimensional
$AB$-image. This flattening has rank at most $n$ if and only if
$E_{AB}(T_j)\subseteq E_{AB}(T_0)$ for every $j$. The other two balanced
flattenings give the corresponding conditions on $AC$ and $BC$.
Therefore $F\in\mathcal F_n$ precisely when every $T_j$ belongs to
$K(T_0)=\C T_0$, which is precisely $F\in X$. Conversely, every
$T\otimes v$ has each balanced rank at most $n$, so $X\subseteq\mathcal F_n$.
This proves~\eqref{eq:components-chart-equality}.

The locus $X$ is the rank-at-most-one variety of $n^3$-by-$m$ matrices.
It is closed by its two-by-two minors and irreducible as the image of
$W\times D$ under $(T,v)\mapsto T\otimes v$. Its dimension is
$n^3+m-1$: on a chart with a chosen nonzero matrix entry, a rank-one
matrix is determined by its corresponding row and column, with their
common entry counted once. If an irreducible closed subset $C$ of
$\mathcal F_n$ contains $X$, then $C\cap U$ is nonempty when
$\Omega_n\ne\varnothing$. It is dense in $C$, and
\eqref{eq:components-chart-equality} places it in the closed subset
$X$. Hence $C=X$, proving the component assertion.

For $v\ne0$, Lemma~\ref{lem:append} gives
$\AR(T\otimes v)=\AR(T)$; the zero tensor is covered by taking $T=0$.
This proves~\eqref{eq:components-sublevel-X}. If $a<n+1$,
Lemma~\ref{lem:flatten} and integrality give
$Z_a(A,B,C,D)\subseteq\mathcal F_n$. The isomorphism in
\eqref{eq:components-sublevel-chart} is
\[
 (T,\lambda_1,\ldots,\lambda_{m-1})
   \longmapsto T\otimes\left(d_0+\sum_{j=1}^{m-1}\lambda_jd_j\right).
\]
Its inverse reads the first slice $T\ne0$ and the unique scalar multiples
of it in the other slices. Those scalars are regular on each chart on
which a selected coordinate of $T$ is nonzero, and these expressions
agree on overlaps. This also proves the claimed isomorphism globally
on the indicated open subset.
\end{proof}

The component assertion concerns the determinantal locus $\mathcal F_n$;
it does not assert that the whole of $X$ lies in an asymptotic-rank
sublevel. Equation~\eqref{eq:components-sublevel-X} retains the actual
three-mode sublevel. The exact sublevel identifications above use four
factors, with $D$ a single factor. When $D$ abbreviates several original
factors below, only the flattening statement
\eqref{eq:components-chart-equality} is used.

\subsection{Scalar-kernel cores in every dimension}
\label{sec:general-cores}

In this subsection $A=B=C=\C^n$, and $K(T)$ is the intersection
of lifted two-mode images defined above. We write $\underline{\R}(T)$
for border rank: the least integer $r$ for which $T$ belongs to the
Zariski closure of the rank-at-most-$r$ locus.

For an integer $n\geq3$ and $s\in\C\setminus\{0\}$, put
\begin{equation}\label{eq:general-core}
 w=\sum_{i=1}^n e_i,
 \qquad T_{n,s}=\sum_{i=1}^n e_i\otimes e_i\otimes e_i
                    +s\,w\otimes w\otimes w.
\end{equation}

\begin{lemma}\label{lem:general-core}
For every $n\geq3$ and every nonzero complex $s$, the tensor
$T_{n,s}$ is concise and $K(T_{n,s})=\C T_{n,s}$.
\end{lemma}

\begin{proof}
Its slices in any mode are $E_{ii}+sJ$, where $J$ is the all-ones
matrix. A linear relation among these slices has the form
\[
 \operatorname{diag}(x_1,\ldots,x_n)
                      +s\Bigl(\sum_i x_i\Bigr)J=0.
\]
An off-diagonal entry gives $\sum_i x_i=0$, and the diagonal entries
then give $x_i=0$ for every $i$. This proves concision in each mode.

Let $S\in K(T_{n,s})$. Every matrix in each two-mode image is
symmetric, so the three lifted-image conditions make $S$ invariant
under all permutations of its indices. In the slice with third
index $k$, write $q_k$ for the common off-diagonal entry. Expressing
this slice in the basis $E_{ii}+sJ$ gives
\[
 q_k=s\left(\sum_i S_{iik}-nq_k\right).
\]
For distinct $j,k$, choose $i$ different from both. Symmetry gives
$S_{ijk}=S_{ikj}$, hence $q_k=q_j$. Let their common value be $q$.
For $i\ne k$, symmetry also gives $S_{iik}=S_{iki}=q$.
Writing $b_k=S_{kkk}$, the displayed relation becomes
$q=s(b_k-q)$, so $b_k=(1+1/s)q$. Thus the coefficients of $S$ are
$q$ off the full diagonal and $(1+1/s)q$ on it, which means
$S=(q/s)T_{n,s}$. The reverse inclusion follows from
$T_{n,s}\in K(T_{n,s})$.
\end{proof}

\begin{lemma}\label{lem:general-core-rank}
For every $n\geq3$ and every nonzero complex $s$,
\begin{equation}\label{eq:general-core-ranks}
 \R(T_{n,s})=\underline{\R}(T_{n,s})=n+1.
\end{equation}
\end{lemma}

\begin{proof}
Equation~\eqref{eq:general-core} gives the upper bound $n+1$ for
both ranks. A concise tensor of rank at most $n$ has rank exactly
$n$, since each single-mode flattening has rank $n$. In an
$n$-term decomposition, its factors span each of the three
$n$-dimensional modes and hence form bases. It is therefore an
invertible local transform of
\[
 D_n=\sum_{i=1}^n e_i\otimes e_i\otimes e_i.
\]
For $D_n$, the lifted-image conditions force all three indices to
coincide, so $K(D_n)$ is the $n$-dimensional span of its diagonal
summands. Invertible local maps carry $K(T)$ to the corresponding
kernel for the transformed tensor. Consequently every concise
tensor of rank at most $n$ has kernel dimension $n$.

On the concise open subset, the appended-minor construction of
Section~\ref{sec:geometry} works with $n^2$-by-$n$ flattenings:
append each column of a variable tensor and impose every
$(n+1)$-row determinant. The resulting polynomial matrix $L_n(T)$
has $n^3$ columns and satisfies $\ker L_n(T)=K(T)$ there.
The condition $\dim K(T)\geq n$ is therefore relatively closed,
being the condition $\rank L_n(T)\leq n^3-n$.
Rank-at-most-$n$ tensors remain dense in their border-rank closure
after restriction to this open subset. Every concise tensor of
border rank at most $n$ must consequently satisfy
$\dim K(T)\geq n$. Lemma~\ref{lem:general-core} excludes this
possibility for $T_{n,s}$, whose kernel has dimension one.
Thus $\underline{\R}(T_{n,s})>n$, proving both equalities.
\end{proof}

\begin{proposition}\label{prop:components-core}
For every integer $n\geq3$, define
\begin{equation}\label{eq:components-Bn}
 B_n=\left((n+1)^n+n^{2n/3}
            -\bigl((n+1)^n-n^n\bigr)^{2/3}\right)^{1/n},
\end{equation}
where the roots are positive real roots. Then $n<B_n<n+1$, and
for every $s\in\C\setminus\{0\}$ the tensor $T_{n,s}$ is concise,
with
\[
 K(T_{n,s})=\C T_{n,s},\qquad
 \R(T_{n,s})=\underline{\R}(T_{n,s})=n+1,\qquad
 \AR(T_{n,s})\leq B_n.
\]
\end{proposition}

\begin{proof}
The concision, kernel and exact-rank assertions are
Lemmas~\ref{lem:general-core} and~\ref{lem:general-core-rank}.
For the asymptotic bound we use the automatic three-direction
form of the speedup theorem of Alman and Li \cite[Theorem~6.1]{AlmanLi2026v2}. For a tensor
$Q$ of format $N\times N\times N$ and an integer ordinary-rank
upper bound $R>2N$, its automatic clause gives
\begin{equation}\label{eq:automatic-speedup}
 \AR(Q)\leq R+N^{2/3}-(R-N)^{2/3}<R.
\end{equation}
Here the auxiliary matrix-rank bound is $N$ in each of the three
directions, which is automatic for $N$-by-$N$ matrices. We apply
this theorem only to three-mode tensors.

Take $Q=T_{n,s}^{\boxtimes n}$, $N=n^n$ and $R=(n+1)^n$.
Taking the grouped $n$th power of the explicit $(n+1)$-term
decomposition~\eqref{eq:general-core} gives $\R(Q)\leq R$. The binomial expansion shows
\[
 (1+1/n)^n
   =2+\sum_{j=2}^n\binom nj n^{-j}>2,
\]
so $R>2N$. Grouped powers satisfy
$\AR(T^{\boxtimes n})=\AR(T)^n$, directly from the defining
limit~\eqref{eq:asymrank}. Substitution in
\eqref{eq:automatic-speedup} yields
$\AR(T_{n,s})^n\leq B_n^n$.

For the endpoints put $z=R-N>N>1$. Since
$N^{2/3}<z^{2/3}$, one has $B_n^n<R$. Also
\[
 B_n^n-N=z-z^{2/3}+N^{2/3}>0.
\]
Taking positive $n$th roots gives $n<B_n<n+1$.
\end{proof}

Thus these cores are available throughout the interval
$B_n\leq a<n+1$: the lower endpoint is included because the
asymptotic-rank estimate is non-strict. The upper endpoint is
excluded because the component argument uses integral flattening
ranks bounded above by $a<n+1$. No exact value of
$\AR(T_{n,s})$ is asserted.
For dimension three, the permutation tensor $P$ supplies the
separate interval $\sqrt[3]{62}\leq a<4$: its scalar kernel is
Lemma~\ref{lem:scalar}, and its strict cost and threshold bounds
are~\eqref{eq:strict-cost} and~\eqref{eq:threshold}.

\subsection{A lower bound for the number of components}
\label{sec:components-count}

\begin{theorem}\label{thm:components-count}
Let $n\geq3$ and $k\geq4$ be integers, let $d_1,\ldots,d_k\geq2$, and put
\[
 t_n=\#\{i:d_i\geq n\}.
\]
For every real $a$ satisfying $B_n\leq a<n+1$, with $B_n$ as in
\eqref{eq:components-Bn}, the variety
$Z_a(\C^{d_1},\ldots,\C^{d_k})$ has at least
\begin{equation}\label{eq:components-count}
                         \binom{t_n}{3}+1
\end{equation}
irreducible components. The same bound with $t_3$ holds for every
$\sqrt[3]{62}\leq a<4$.
\end{theorem}

\begin{proof}
First suppose that every $d_i\leq n$, and call a factor active if its
dimension is $n$. For the first interval, choose any nonzero $s$ and
put $Q=T_{n,s}$ from Proposition~\ref{prop:components-core}. For the
second interval put $n=3$ and $Q=P$. Lemma~\ref{lem:scalar} and
\eqref{eq:strict-cost} supply its scalar kernel and cost bound. Thus
in both cases $Q$ is concise, $K(Q)=\C Q$, and
\[
                         \AR(Q)\leq a<n+1.
\]
In particular $a\geq n\geq3$ by Lemma~\ref{lem:flatten}.

For each triple $I$ of active factors, let $x_I$ be $Q$ on those
factors and a fixed basis vector $e_0$ on every other factor.
Repeated application of Lemma~\ref{lem:append} gives
$\AR(x_I)=\AR(Q)$, so $x_I\in Z_a$.
For linear-algebra purposes write
$D=\bigotimes_{j\notin I}\C^{d_j}$ and choose its product basis.
The complementary all-zero coordinate selects the first slice $Q$.
Proposition~\ref{prop:components-chart} supplies a full-ambient open
neighborhood $U_I$ of $x_I$. Every bipartition rank on $Z_a$ is an
integer at most $n$, so
\begin{equation}\label{eq:components-local}
 Z_a\cap U_I\subseteq X_I,
 \qquad X_I=\{F:\rank F_{I\mid I^c}\leq1\}.
\end{equation}
This use of $D$ does not merge any modes in the definition of $Z_a$.

The closed affine variety $Z_a$ has finitely many irreducible components.
Choose a component $C_I$ containing $x_I$. Its nonempty relatively
open subset $C_I\cap U_I$ is dense in $C_I$ and lies in the closed
locus $X_I$. It follows that
\begin{equation}\label{eq:components-containment}
                              C_I\subseteq X_I.
\end{equation}
No uniqueness of a component through $x_I$ is needed.

There is a useful stronger constraint: every single-mode rank outside
$I$ is at most one everywhere on $C_I$. Fix $i\in I$. The locus on
$C_I$ where the $i$th single-mode rank is $n$ is nonempty, since it
contains $x_I$, and is open, since that mode has dimension $n$.
At any point $F$ of this open set, \eqref{eq:components-containment}
gives a nonzero factorization $F=A_I\otimes B_{I^c}$. For $j\notin I$,
the Kronecker-product rank formula gives
\[
 \rank F_{\{i,j\}\mid\{i,j\}^c}
      =r_i(A_I)r_j(B_{I^c})=n\,r_j(B_{I^c})\leq n.
\]
Here $r_j$ denotes the appropriate single-mode flattening rank, and
the last inequality follows from $F\in Z_a$. Thus $r_j(F)=1$ on
that open set. The condition $r_j(F)\leq1$ is closed, so density
extends it to all of $C_I$. If $I\ne J$, choose $j\in J\setminus I$.
Then $x_J$ has $j$th rank $n$, whereas every point of $C_I$ has
$j$th rank at most one. Therefore the chosen components $C_I$ are
pairwise distinct, whether or not the triples overlap.

Every factor has dimension at least two. The tensor
\[
 g=\bigotimes_{i=1}^k e_0+\bigotimes_{i=1}^k e_1
\]
has rank and asymptotic rank two: its displayed decomposition gives
the upper bound, and every nontrivial bipartition has matrix rank two.
Thus $g\in Z_a$. Since $k\geq4$, each $I\mid I^c$ is nontrivial,
so $g\notin X_I$ for every active triple. A component containing $g$
is different from all the $C_I$, proving~\eqref{eq:components-count}
in these dimensions. If there are no active triples, the bound is
one and follows from the same nonempty sublevel.

For arbitrary $d_i\geq2$, choose coordinate projections
\[
 p_i:\C^{d_i}\longrightarrow\C^{\min(d_i,n)}
\]
with coordinate inclusions $j_i$ as sections. The tensor-product maps
$p=\bigotimes_i p_i$ and $j=\bigotimes_i j_i$ are local maps for the
original $k$ factors and satisfy $pj=1$. Rank is nonincreasing under
local maps in every power. Consequently $p$ maps the larger sublevel
onto the smaller one, and $j$ preserves asymptotic rank on the embedded
copy.

If the larger sublevel has components $C_1,\ldots,C_s$, the closures
of $p(C_1),\ldots,p(C_s)$ in the smaller sublevel are irreducible
closed sets covering it. Any irreducible component of the smaller
sublevel lies in one of those finitely many closed sets, and maximality
makes it equal to that set. The smaller sublevel therefore has at most
$s$ components. Its active-factor count is exactly $t_n$, so its
established lower bound gives the same bound in the original dimensions.
\end{proof}

When all $d_i\geq n$, this lower bound is $\binom{k}{3}+1$.
In particular, four factors of dimensions at least three have at least
five components for every $\sqrt[3]{62}\leq a<4$. The counts are lower
bounds, not complete lists of components. Their unboundedness varies
the tensor format; it makes no claim of infinitely many asymptotic-rank
values in a fixed format. The same section-and-retraction argument
also propagates reducibility when additional factors are appended
using fixed nonzero vectors.

\subsection{The grouped three-mode sublevel}
\label{sec:components-grouping}

\begin{proposition}\label{prop:components-grouping}
For every real $3\leq a<4$,
\begin{equation}\label{eq:components-grouped}
 Z_a(\C^9,\C^3,\C^2)
   =\{F:\rank F_{\C^9\mid\C^3\otimes\C^2}\leq3\}.
\end{equation}
This is the irreducible rank-at-most-three determinantal variety of
$9$-by-$6$ matrices.
\end{proposition}

\begin{proof}
Lemma~\ref{lem:flatten} gives the forward inclusion, using integrality
and $a<4$. Conversely, a tensor with this flattening rank at most three
has its first-mode support in a subspace of dimension at most three.
After compression and padding within that mode, it suffices to consider
format $(3,3,2)$.

Represent such a tensor by a pair of three-by-three matrices $(A,B)$.
The set where $A$ is invertible and $A^{-1}B$ has three distinct
eigenvalues is nonempty, as shown by $(I,\operatorname{diag}(0,1,2))$.
It is Zariski open: on $\det A\ne0$, distinct eigenvalues are the
nonvanishing of the characteristic discriminant. Thus it is dense in
the affine space of all pairs. Invertible left and right changes of
basis reduce every such pair to
\[
                  (I,\operatorname{diag}(\lambda_1,\lambda_2,\lambda_3)).
\]
The corresponding tensor is a sum of the three decomposable tensors
$e_i\otimes e_i\otimes(d_0+\lambda_i d_1)$. Its rank is at most three.
Taking Zariski closures proves that every tensor of format $(3,3,2)$
has border rank at most three and hence asymptotic rank at most three.
The latter implication also follows directly from closedness of $Z_3$:
that closed set contains every rank-at-most-three tensor and therefore
its closure. Local padding back to $(9,3,2)$ cannot increase the
asymptotic rank. Since $a\geq3$, this proves the reverse inclusion.

Every $9$-by-$6$ matrix of rank at most three is a product of a
$9$-by-$3$ matrix and a $3$-by-$6$ matrix. Hence the determinantal
variety is the image of an irreducible affine space under matrix
multiplication and is irreducible.
\end{proof}

Grouping any two of the three-dimensional modes of $(3,3,3,2)$ gives
$(9,3,2)$. Throughout $\sqrt[3]{62}\leq a<4$, the full four-mode
sublevel is reducible by Theorem~\ref{thm:components-count}, whereas
the full grouped three-mode sublevel is irreducible by
Proposition~\ref{prop:components-grouping}. This does not assert
surjectivity from the ungrouped sublevel onto the grouped one and
does not settle the general three-mode irreducibility question.

\subsection{Scope and the remaining three-mode questions}
\label{sec:components-scope}

The component lower bounds concern the full sublevels and allow the
format to vary. They neither classify all components nor produce
infinitely many asymptotic-rank values in a single fixed format.
The scalar-kernel argument uses three balanced flattenings and a
fourth factor containing simultaneous slices. The grouped example
shows explicitly why those constraints do not transfer to three
modes. A reducible linear section would also be insufficient to
prove reducibility of an ambient sublevel: the irreducible quadric
$xy=z^2$ has the reducible section $z=0$.

The full three-mode asymptotic-rank conjecture asserts
$\AR(T)=\max_i r_i(T)$, where $r_i(T)$ are the single-mode ranks,
after compression to those mode supports. Under this conjecture,
for every $a\geq0$ the sublevel is the subspace variety defined by
\[
 r_i(T)\leq\min\{\dim V_i,\lfloor a\rfloor\},\qquad i=1,2,3.
\]
It is irreducible: choose the three support spaces in the product
of the corresponding Grassmannians and then a tensor in their
tensor product. The resulting vector bundle has irreducible total
space and maps onto the subspace variety. When the bound is zero,
the variety is $\{0\}$. This conditional implication is already
recorded in \cite[Section~5]{CHNVZ2024}. Thus three-mode
reducibility would refute the full conjecture. It would not by
itself prove that the matrix-multiplication exponent $\omega$ is
greater than two, or refute the narrower conjecture for concise
tight cubic tensors \cite[Conjectures~3 and~8]{BjorklundKaski2023}.

\begin{proposition}\label{prop:components-finite-values}
In a fixed complex tensor format, finiteness of the set of
asymptotic-rank values is equivalent to a uniform bound on the
number of irreducible components of all its sublevels.
\end{proposition}

\begin{proof}
Finitely many values give finitely many sublevels, each a closed
affine variety with finitely many irreducible components.
Conversely, suppose all sublevels in an ambient affine space of
dimension $N$ have at most $C$ components. Associate to each one
its vector of component counts in dimensions $N,N-1,\ldots,0$.
For a strict inclusion of two closed subsets these vectors
strictly increase lexicographically. Indeed, until the first
dimension at which their components differ, all higher-dimensional
components coincide. An old component in the next dimension
cannot be absorbed by one of those common higher-dimensional
components, and containment in a new component of equal dimension
forces equality. Thus all old components of that dimension persist,
and at least one new component is added. If there were no differing
dimension, the two closed subsets would be equal.

There are only $\binom{C+N+1}{N+1}$ nonnegative integer vectors of
length $N+1$ with coordinate sum at most $C$. Consequently every
strict chain of sublevels has uniformly bounded finite length.
Distinct attained asymptotic-rank values give distinct, totally
ordered sublevels, so there are only finitely many such values.
\end{proof}

No uniform component bound in a fixed format is supplied by the
component-growth theorem. Determining the three-mode components in
the local product chart, finding a three-mode replacement for the
balanced-flattening argument, and controlling components over all
thresholds of a fixed format remain separate questions.

The arguments in this section are written proofs with the cited
closedness and asymptotic-speedup inputs. The finite certificate
described in Section~\ref{sec:verification} supports the original
permutation-tensor construction. It does not formalize the general
scalar cores, the component-growth theorem, the grouped
determinantal description, or the fixed-format proposition above.
